\section{Introduction}

The aim of this project was to implement a small bare-metal operating system for the Raspberry Pi4. The system is loaded directly as a kernel image and does not rely on Linux, a standard C runtime, processes, filesystems, or user-space services. Its main purpose is to demonstrate how core operating system mechanisms can be built from scratch on real hardware.

The project focusses on bootstrapping, kernel initialization, cooperative multitasking, context switching, synchronization, hardware drivers, and interactive system control. The resulting system is not intended to be a general-purpose operating system. Instead, it is a compact embedded kernel that can schedule multiple kernel tasks, communicate over UART, render text output to HDMI, access external hardware over I2C, and provide runtime diagnostics.